% ===================================================================
%  TABLEAU N° 10 — La mer. Le port.
%
%  Traduction, faite en 2026, en francais standard du Quebec, sur
%  l'ido de texto/io. Regles de la colonne : voir 10-tableau-01.tex.
%
%  LE TABLEAU LE MOINS QUEBECOIS DE LA SERIE, ET IL FAUT LE DIRE
%  PLUTOT QUE DE MEUBLER LA LISTE. Trois choix seulement :
%    (42) cargo-boat .......... cargo
%    (75) halte ............... arrêt
%    (3)  bateaux à vapeur .... navires à vapeur
%
%  LA RAISON EST DANS LA MATIERE. Le vocabulaire de la marine est un
%  vocabulaire de METIER, fixe par la meme ecole et les memes manuels
%  des deux cotes de l'Atlantique : la quille, le gouvernail, la
%  vergue, le chenal, la bouee, l'ecluse, le sémaphore, le
%  quartier-maitre se disent ici comme la-bas, et se disaient deja
%  ainsi en 1926. Il n'y avait rien a moderniser : les rares mots
%  vieillis de ce tableau — gabare, chaland, locomobile — nomment des
%  bateaux et des machines qui ont disparu avec leur nom. On ne
%  remplace pas le nom d'une chose qui n'existe plus : ce serait
%  moderniser la chose, non la langue.
%
%  « JETEE » RESTE, ET C'EST UN CAS CURIEUX. L'ido ecrit « warfo »,
%  qui est l'anglais \textit{wharf}, et le mot que le Quebec emploierait
%  spontanement pour cet ouvrage est « quai ». Mais « quai » traduit
%  deja « kayo » (58), qui est un autre ouvrage du meme port et porte
%  son propre renvoi. Les deux ne peuvent pas partager un mot.
% ===================================================================

%%K t10-apar-1 apar
\VUsaut{4.0mm}
\VUtitre{15.0pt}{60}{TABLEAU N\textsuperscript{o} 10}
\VUsaut{3.0mm}
\VUpk{11.4pt}{40}{La mer. --- Le port.}
\VUsaut{3.0mm}

%%K t10-01-1 p
1. --- L'été, pendant que les uns cherchent le calme et la fraîcheur à la
montagne, les autres aiment mieux se rendre au bord de la mer. C'est ce que nous
avons fait l'an dernier, parce que nous aimons beaucoup l'\VUgras{Océan}
\textsuperscript{(1)}.

%%K t10-02-1 p
2. --- Pour ma part, j'éprouve un grand plaisir à contempler cette immense
étendue d'eau qui va jusqu'à l'\VUgras{horizon} \textsuperscript{(2)}, et à voir
partir pour une longue \VUgras{traversée} les énormes \VUgras{navires à vapeur}
\textsuperscript{(3)} sur lesquels s'embarquent les voyageurs qui vont en
Amérique.

%%K t10-03-1 p
3. --- Aussi mon père, qui connaît mon penchant, choisit-il toujours, pour
passer les vacances, une plage bien tranquille, située près d'un de nos grands
\VUgras{ports militaires}.

%%K t10-03-2 p
Pendant notre dernier séjour, l'\VUgras{escadre} est venue jeter l'ancre
derrière l'énorme \VUgras{digue} \textsuperscript{(4)} qui ferme et protège le
\VUgras{port militaire} \textsuperscript{(5)}, et j'ai pu admirer à mon aise les
différents \VUgras{navires de guerre}, depuis le petit \VUgras{sous-marin}
\textsuperscript{(10)}, qui plonge et disparaît sous les eaux, jusqu'à l'énorme
\VUgras{cuirassé} \textsuperscript{(9)} qui, jusqu'ici, ne craignait guère que
les attaques du \VUgras{torpilleur} \textsuperscript{(8)}.

%%K t10-tit-2 sub
\VUsaut{3.0mm}
\VUpk{10.2pt}{40}{Les petits navires.}
\VUsaut{3.0mm}

%%K t10-04-1 p
4. --- Celui-ci savait déjà trouver très \VUgras{habilement} le point faible de
l'épaisse cuirasse de métal pour y planter la dangereuse \VUgras{torpille},
dont l'explosion cause la perte du navire; mais il était pourtant moins à
craindre que le sous-marin, parce que les contre-torpilleurs le poursuivent
facilement.

%%K t10-05-1 p
5. --- J'ai pu voir aussi les rapides \VUgras{croiseurs} \textsuperscript{(6)}
cuirassés, presque aussi invulnérables que le vrai cuirassé, plusieurs
\VUgras{transports} \textsuperscript{(12)}, trois ou quatre \VUgras{canonnières}
\textsuperscript{(7)} et enfin une \VUgras{frégate} \textsuperscript{(11)}.
\VUgras{Le spectacle de cette flotte, quand elle manœuvre pour lever l'ancre,
est des plus intéressants.}

%%K t10-06-1 p
6. --- Quelle différence de construction entre ces grandes masses d'acier et les
élégants \VUgras{trois-mâts} ou les fins \VUgras{voiliers} \textsuperscript{(13)}
qu'on emploie encore aujourd'hui dans la flotte de commerce, et que j'ai vus
entrer au port, toutes voiles dehors, pendant que je me promenais sur la
\VUgras{jetée} \textsuperscript{(14)}. Il est vrai que ceux-ci perdaient
beaucoup de leurs avantages \VUgras{quand} le \VUgras{vent} était contraire. Ils
devaient amener toutes leurs voiles pour rentrer dans le bassin, et il fallait
un \VUgras{remorqueur} \textsuperscript{(16)} pour aller les chercher et les
amener au quai comme de simples \VUgras{chalands} \textsuperscript{(17)}.

%%K t10-tit-3 sub
\VUsaut{3.0mm}
\VUpk{10.2pt}{40}{Le port de commerce.}
\VUsaut{3.0mm}

%%K t10-07-1 p
7. --- Ce qu'il y a d'intéressant dans le port dont je parle, c'est qu'il
contient non seulement un port militaire, aménagé artificiellement par des
travaux gigantesques, mais aussi un merveilleux \VUgras{port de commerce} situé
à l'\VUgras{embouchure} d'un grand fleuve.

%%K t10-08-1 p
8. --- La langue de terre qui sépare les deux bassins est fortifiée et défendue
par une série de \VUgras{batteries} et de \VUgras{bastions} qui s'avancent sur
la mer. Il faut une permission spéciale pour se promener sur les
\VUgras{remparts} \textsuperscript{(15)}. Je l'avais obtenue facilement par mon
oncle, qui est ingénieur des \VUgras{chantiers de construction}
\textsuperscript{(18)}, et je suis allé visiter les navires qu'on construit sous
de vastes hangars.

%%K t10-09-1 p
9. --- J'ai même assisté au lancement d'un cuirassé, et je me rappelle l'instant
émouvant que toute la foule a ressenti quand l'énorme masse, abandonnée à
elle-même, s'est mise à glisser sur son berceau et est venue flotter sur l'eau
du fleuve.

%%K t10-10-1 p
10. --- Pour se rendre aux chantiers, on traverse le \VUgras{terrain de
manœuvres} \textsuperscript{(19)} où les soldats de la garnison viennent
s'exercer chaque jour, et l'on passe sur une \VUgras{passerelle}
\textsuperscript{(20)} de fer qui relie l'esplanade de parade à
l'\VUgras{arsenal} \textsuperscript{(21)}.

%%K t10-tit-4 sub
\VUsaut{3.0mm}
\VUpk{10.2pt}{40}{L'arsenal et les bassins.}
\VUsaut{3.0mm}

%%K t10-11-1 p
11. --- Avec mon oncle, j'ai visité ce grand établissement plein de
\VUgras{canons} \textsuperscript{(22)}, de \VUgras{boulets}
\textsuperscript{(23)} et d'\VUgras{obus}. J'ai vu aussi l'\VUgras{usine
métallurgique} \textsuperscript{(24)} où l'on fabrique les \VUgras{chaudières}
\textsuperscript{(25)} des navires à vapeur.

%%K t10-12-1 p
12. --- Tout près de là se trouvent les \VUgras{bassins de radoub}
\textsuperscript{(26)} et les très curieux \VUgras{docks flottants}
\textsuperscript{(27)}, où j'ai pu voir de tout près, sur un navire en
réparation, l'\VUgras{hélice} \textsuperscript{(28)}, la \VUgras{quille}
\textsuperscript{(29)} et le \VUgras{gouvernail} \textsuperscript{(30)}.

%%K t10-13-1 p
13. --- Le port de commerce est tout aussi digne d'étude que le port militaire,
et j'ai passé bien des heures à flâner sur les quais où sont amarrés les
\VUgras{bateaux à aubes} \textsuperscript{(31)} qui remontent le fleuve, et à
regarder fonctionner les \VUgras{grues à mâter} \textsuperscript{(32)}, qui
soulèvent comme des plumes d'énormes fardeaux.

%%K t10-tit-5 sub
\VUsaut{4.0mm}
\VUpk{10.2pt}{40}{Le pilote.}
\VUsaut{3.0mm}

%%K t10-14-1 p
14. --- J'admirais les \VUgras{transatlantiques} \textsuperscript{(34)} sortant
de la rade, \VUgras{pavillons} \textsuperscript{(35)} au vent, conduits par un
habile \VUgras{pilote} \textsuperscript{(36)} qui sait merveilleusement suivre
le \VUgras{chenal} marqué par les \VUgras{bouées} \textsuperscript{(40)} et les
\VUgras{balises}, en évitant les \VUgras{gabares} \textsuperscript{(41)} qui se
dirigent péniblement avec leur unique voile carrée. Je distinguais à
l'\VUgras{avant} \textsuperscript{(37)} --- la proue --- les \VUgras{cabines}
\textsuperscript{(39)} de deuxième classe et, à l'\VUgras{arrière}
\textsuperscript{(38)} --- la poupe --- les salons des passagers de première
classe.

%%K t10-15-1 p
15. --- Parfois, j'assistais aussi au chargement d'un \VUgras{cargo}
\textsuperscript{(42)} dans lequel on venait d'entasser les marchandises de
l'\VUgras{entrepôt} \textsuperscript{(43)} et de la \VUgras{gare maritime}
\textsuperscript{(44)}, amenées par les nombreux trains de la \VUgras{ligne des
quais} \textsuperscript{(45)}.

%%K t10-tit-6 sub
\VUsaut{3.0mm}
\VUpk{10.2pt}{40}{Le canotage.}
\VUsaut{3.0mm}

%%K t10-16-1 p
16. --- J'ai souvent fait du canot dans le \VUgras{bassin à flot}
\textsuperscript{(53)}, où viennent se réfugier les \VUgras{embarcations}
\textsuperscript{(46)}, et c'était une joie pour moi de manier l'\VUgras{aviron}
\textsuperscript{(47)}. Je montais sur le \VUgras{pont} \textsuperscript{(48)}
d'une \VUgras{barque à voile} \textsuperscript{(49)}, je grimpais par les
\VUgras{cordages} \textsuperscript{(52)} le long du \VUgras{mât}
\textsuperscript{(51)} jusque dans les \VUgras{vergues} \textsuperscript{(50)},
puis je reprenais ma promenade en \VUgras{yole} \textsuperscript{(54)} entre les
lourdes \VUgras{barques de pêcheurs} \textsuperscript{(55)}, où sèchent les
\VUgras{filets} \textsuperscript{(56)}.

%%K t10-17-1 p
17. --- Sur le \VUgras{quai} \textsuperscript{(58)} défilaient sous mes yeux les
\VUgras{marchandises} \textsuperscript{(59)} les plus diverses, contenues dans
des \VUgras{caisses} \textsuperscript{(60)} ou dans des \VUgras{ballots}
\textsuperscript{(61)}. Elles formaient la \VUgras{cargaison}
\textsuperscript{(63)} des \VUgras{péniches}, et de puissantes \VUgras{grues}
\textsuperscript{(62)} les enlevaient et les déposaient sur des
\VUgras{camions} \textsuperscript{(64)} pendant que les \VUgras{camionneurs} les
déclaraient et payaient les droits au \VUgras{bureau de la douane}
\textsuperscript{(65)}.

%%K t10-18-1 p
18. --- Enfin, quand j'arrivais au \VUgras{pont tournant} \textsuperscript{(66)}
ou à l'\VUgras{écluse} \textsuperscript{(69)}, en passant devant la
\VUgras{drague} \textsuperscript{(68)} qui nettoie le \VUgras{canal}
\textsuperscript{(67)}, je débarquais sur la jetée, à côté du
\VUgras{sémaphore} \textsuperscript{(70)}.

%%K t10-19-1 p
19. --- C'est de l'autre côté de cette jetée que viennent accoster les
\VUgras{chaloupes} \textsuperscript{(71)} qui amènent à terre les officiers de
l'escadre. C'est un spectacle admirable de voir les matelots lever en cadence
leurs avirons, pendant que l'embarcation, portée par la \VUgras{vague}
\textsuperscript{(72)}, aborde facilement l'escalier du débarcadère. C'est là
qu'on peut le mieux observer la \VUgras{marée} et le mouvement du \VUgras{flux}
et du \VUgras{reflux} sur la \VUgras{plage} \textsuperscript{(73)}.

%%K t10-tit-7 sub
\VUsaut{3.0mm}
\VUpk{10.2pt}{40}{Le cercle.}
\VUsaut{3.0mm}

%%K t10-20-1 p
20. --- Pour rentrer à la maison, je prenais le \VUgras{tramway électrique}
\textsuperscript{(74)}, dont l'\VUgras{arrêt} \textsuperscript{(75)} est près du
\VUgras{poteau} placé devant le cercle des officiers.

%%K t10-20-2 p
Du \VUgras{balcon} \textsuperscript{(76)} du cercle, décoré d'\VUgras{ancres}
\textsuperscript{(77)}, de canons et de boulets, j'ai souvent vu une splendide
assemblée d'officiers de marine : des \VUgras{lieutenants de vaisseau}
\textsuperscript{(78)} avec leur épée, des \VUgras{capitaines d'artillerie}
\textsuperscript{(79)} et d'\VUgras{infanterie} \textsuperscript{(80)} avec leur
\VUgras{sabre}. Derrière eux se tenait un \VUgras{matelot}
\textsuperscript{(81)} qui leur apportait une \VUgras{longue-vue} quand ils
voulaient distinguer ce qui se passait au loin.

%%K t10-21-1 p
21. --- J'ai vu aussi de jeunes \VUgras{enseignes} \textsuperscript{(82)}
recevant les instructions de leur \VUgras{capitaine de vaisseau}
\textsuperscript{(83)} ou de l'\VUgras{amiral} \textsuperscript{(84)}, pendant
qu'un \VUgras{quartier-maître} \textsuperscript{(85)} attendait pour les porter
au navire.

%%K t10-22-1 p
22. --- De ce balcon, la vue est splendide : on domine toute la plage avec ses
\VUgras{tentes} \textsuperscript{(86)} et ses \VUgras{cabines}
\textsuperscript{(87)}, et l'on voit, à l'heure du bain, les \VUgras{baigneurs}
\textsuperscript{(88)} qui s'ébattent joyeusement devant le \VUgras{casino}
\textsuperscript{(89)}.

%%K t10-23-1 p
23. --- Au bout de la plage, la côte forme une petite \VUgras{presqu'île}
\textsuperscript{(91)} de \VUgras{rochers}, où la \VUgras{lame}
\textsuperscript{(92)} vient se briser et sur laquelle est bâti un
\VUgras{phare} \textsuperscript{(93)} qui indique la nuit la route aux
navigateurs. Cette presqu'île n'est reliée à la terre que par un \VUgras{isthme}
très étroit \textsuperscript{(90)}.

%%K t10-tit-8 sub
\VUsaut{3.0mm}
\VUpk{10.2pt}{40}{Vue générale de la côte.}
\VUsaut{3.0mm}

%%K t10-24-1 p
24. --- La \VUgras{falaise} \textsuperscript{(94)} s'étend, découpée par la mer,
qui y dessine capricieusement une série de \VUgras{baies} \textsuperscript{(95)}
et de \VUgras{promontoires} (des caps), ce qui la rend des plus pittoresques.

%%K t10-24-2 p
Sur le \VUgras{plateau} \textsuperscript{(96)}, qui domine le \VUgras{détroit}
\textsuperscript{(98)}, il y a un \VUgras{fort} \textsuperscript{(97)} très
important et quelques « villas ».

%%K t10-25-1 p
25. --- Ce détroit sépare du continent une \VUgras{île} \textsuperscript{(99)}
très dangereuse, bordée d'\VUgras{écueils} \textsuperscript{(100)} et de
\VUgras{brisants}, où des barques et même de grands navires s'échouent parfois.
Malgré la promptitude des secours et l'arrivée rapide des \VUgras{canots de
sauvetage}, ces \VUgras{naufrages} sont terrifiants et, pendant les
\VUgras{tempêtes} d'équinoxe, plusieurs navires ont péri corps et biens.

%%K t10-25-2 p
Malgré ce voisinage, \VUgras{le port} est très fréquenté par les marins.
\VUsaut{6.0mm}
\VUfilet{22mm}
